\begingroup
\setlength{\abovedisplayskip}{7pt}
\setlength{\belowdisplayskip}{7pt}
\setlength{\abovedisplayshortskip}{3pt}
\setlength{\belowdisplayshortskip}{4pt}

% Printed 52 (unfoliated); principal PDF 59; Google PDF 63.
\par\clearpage

\vspace*{.19\textheight}
\begin{center}
{\large IV.}\par\medskip
{\Large BREITENBESTIMMUNG\\DER STERNWARTE SEEBERG.}
\end{center}

Gleichzeitig mit meinen Beobachtungen in Göttingen und Altona wurden
dieselben Sterne auf meine Aufforderung auch von Hrn. \textsc{Hansen}, Director der
Sternwarte Seeberg bei Gotha, an dem dortigen \textsc{Ertel}schen zweifussigen
Meridiankreise beobachtet. Der sich daraus ergebende Breitenunterschied zwischen
dieser und der Göttinger Sternwarte erhält ein noch erhöhtes Interesse durch
den Umstand, dass erstere vermittelst einiger unter Leitung des Herrn
Generallieutenants \textsc{von Müffling} gemessener Dreiecke mit dem hannoverschen
Dreieckssystem verbunden ist.

Der Kreis wurde während der Beobachtungen einigemale umgelegt, allein
die Bestimmung des Collimationsfehlers wurde unabhängig davon jeden Tag,
und meistens jeden Tag zweimal, durch Einstellung auf den Nadirpunkt gemacht,
welches schon oben erwähnte Verfahren Hr. \textsc{Hansen} im Herbst 1826
auf hiesiger Sternwarte praktisch kennen gelernt hatte. Die Ablesung geschah
nicht mit Verniers, sondern mit Mikroskopen. Folgende Übersicht enthält
die Hauptresultate dieser Beobachtungen, indem die erste Columne die
Bezeichnung des Sterns, die zweite die Lage des Kreises, die dritte die Anzahl
der Beobachtungen, die vierte die von mir auf den Anfang des Jahrs 1827
reducirte Zenithdistanz (nördliche mit positivem Zeichen), die fünfte die Breite,
welche aus den oben S.~9 mitgetheilten Declinationen sich ergibt, darstellt.

% Printed 53; principal PDF 60; Google PDF 64.
\par\clearpage
\noindent\makebox[12mm][l]{}\parbox[t]{\dimexpr\linewidth-24mm\relax}{\centering\scriptsize\bfseries BESTIMMUNG DES BREITENUNTERSCHIEDES ETC.\quad IV. BREITENBESTIMMUNG VON SEEBERG.}\makebox[12mm][r]{\fontsize{13}{16}\selectfont 53}\par\medskip

{\small
\setlength{\arraycolsep}{3.2pt}
\renewcommand{\arraystretch}{1.06}
\[
\begin{array}{@{}r|c|c|r@{\;}r@{\;}r|r||r|c|c|r@{\;}r@{\;}r|r@{}}
1&\text{Ost}&5&-1^\circ&1'&53''{,}14&50^\circ56'4''{,}76&
16&\text{Ost}&5&-1^\circ&42'&19''{,}87&50^\circ56'6''{,}55\\
&\text{West}&4&&&52{,}80&4{,}42&&\text{West}&5&&&17{,}05&3{,}73\\
2&\text{Ost}&5&+4&37&29{,}41&5{,}57&17&\text{Ost}&5&+0&38&48{,}22&5{,}50\\
&\text{West}&6&&&30{,}86&4{,}12&&\text{West}&5&&&48{,}79&4{,}93\\
3&\text{Ost}&5&-0&45&19{,}00&5{,}20&18&\text{Ost}&5&+1&38&59{,}63&6{,}24\\
&\text{West}&7&&&17{,}70&3{,}90&&\text{West}&6&&&61{,}51&4{,}36\\
4&\text{Ost}&6&+3&38&50{,}38&5{,}17&19&\text{West}&2&+3&41&30{,}38&4{,}49\\
&\text{West}&6&&&51{,}05&4{,}50&20&\text{Ost}&1&+3&33&50{,}19&4{,}11\\
5&\text{Ost}&6&+4&29&54{,}35&4{,}87&&\text{West}&1&&&47{,}87&6{,}43\\
&\text{West}&6&&&54{,}30&4{,}92&21&\text{Ost}&1&+1&58&31{,}74&5{,}35\\
6&\text{Ost}&6&-4&20&26{,}87&6{,}01&&\text{West}&1&&&32{,}90&4{,}19\\
&\text{West}&4&&&25{,}92&5{,}06&22&\text{Ost}&1&-4&39&61{,}37&6{,}07\\
7&\text{Ost}&6&-0&39&22{,}12&5{,}56&&\text{West}&1&&&58{,}90&3{,}60\\
&\text{West}&4&&&20{,}70&4{,}14&23&\text{Ost}&1&+5&24&21{,}09&5{,}91\\
8&\text{Ost}&5&+1&40&1{,}41&6{,}06&&\text{West}&1&&&22{,}88&4{,}12\\
&\text{West}&4&&&3{,}24&4{,}23&24&\text{West}&1&-0&33&26{,}54&5{,}95\\
9&\text{Ost}&5&+5&17&29{,}69&6{,}82&25&\text{West}&1&+8&5&40{,}64&6{,}50\\
&\text{West}&4&&&33{,}72&2{,}79&26&\text{West}&1&-0&17&51{,}91&5{,}03\\
10&\text{Ost}&5&+1&43&6{,}51&5{,}54&27&\text{West}&1&-4&35&46{,}53&3{,}57\\
&\text{West}&5&&&8{,}63&3{,}42&28&\text{West}&1&+2&44&9{,}48&3{,}93\\
11&\text{Ost}&5&+2&43&27{,}76&5{,}26&29&\text{West}&1&+1&31&5{,}78&5{,}28\\
&\text{West}&4&&&28{,}89&4{,}13&30&\text{West}&2&+4&40&0{,}22&4{,}51\\
12&\text{Ost}&5&+2&2&48{,}89&6{,}77&31&\text{West}&1&-4&57&58{,}28&4{,}13\\
&\text{West}&3&&&50{,}29&5{,}37&32&\text{West}&1&+2&18&59{,}76&3{,}88\\
13&\text{Ost}&5&-1&29&56{,}73&6{,}19&33&\text{West}&1&-0&39&51{,}61&5{,}28\\
&\text{West}&5&&&56{,}64&6{,}10&34&\text{West}&1&+6&9&34{,}85&2{,}94\\
14&\text{Ost}&5&-0&35&42{,}51&5{,}00&35&\text{West}&1&-3&59&15{,}56&4{,}09\\
&\text{West}&3&&&42{,}84&5{,}33&36&\text{West}&2&-4&6&42{,}83&4{,}84\\
15&\text{Ost}&5&-2&36&14{,}46&6{,}93&37&\text{West}&1&+6&0&36{,}97&5{,}38\\
&\text{West}&4&&&12{,}88&5{,}35&38&\text{West}&1&-1&53&16{,}75&3{,}84
\end{array}
\]
\par
}

Diese sechzig Resultate für die Breite haben nun freilich ungleiche
Zuverlässigkeit; allein um die ihnen beizulegenden Gewichte ohne Willkür
angeben zu können, müsste das Verhältniss des mittlern eigentlichen
Beobachtungsfehlers zum mittlern Theilungsfehler bekannt sein; ist dies Verhältniss
wie \(1\) zu \(\surd\theta\), so wird, wenn man die geringe den Declinationen noch
anhängende Unsicherheit nicht beachtet,
\[
\frac n{1+n\theta}
\]
\noindent das Gewicht einer auf \(n\) Beobachtungen, die sich auf einerlei Theilstrich
beziehen, beruhenden Bestimmung sein.

Nimmt man statt dieses Gewichts schlechthin \(n\) an, so wird das Mittel
aus den 206 Beobachtungen
\[
=50^\circ56'5''{,}16.
\]

% Printed 54; principal PDF 61; Google PDF 65.
\par\clearpage
\noindent\makebox[12mm][l]{\fontsize{13}{16}\selectfont 54}\parbox[t]{\dimexpr\linewidth-24mm\relax}{\centering\scriptsize\bfseries BESTIMMUNG DES BREITENUNTERSCHIEDES ZWISCHEN GÖTTINGEN UND ALTONA.}\makebox[12mm][r]{}\par\medskip

Inzwischen lassen die Beobachtungen erkennen, dass die Theilungsfehler
vergleichungsweise beträchtlich grösser sein müssen, als an dem \textsc{Ramsden}schen
Zenithsector, während die eigentlichen Beobachtungsfehler eher noch etwas
kleiner sein mögen. Bei jenem Verfahren werden also die auf einer grössern
Anzahl von Beobachtungen beruhenden Bestimmungen vor denen, welchen nur
eine oder zwei zum Grunde liegen, viel zu sehr bevorzugt.

Sobald man aber den Einfluss der Theilungsfehler berücksichtigen will,
darf auch nicht unbeachtet bleiben, dass die jedesmalige Bestimmung des
Collimationsfehlers einen constanten, von den Fehlern der dabei sprechenden
Theilstriche abhängenden Fehler involvirt. Es ist aber klar, dass derselbe auf
die Polhöhe in entgegengesetztem Sinn wirkt, je nachdem der Kreis östlich
oder westlich sich befindet. Man wird daher die auf die verschiedenen Lagen
des Kreises sich beziehenden Beobachtungen von einander trennen, aus jeder
Reihe, mit Anwendung des Gewichts \(\dfrac n{1+n\theta}\) für jede Bestimmung, das Mittel
berechnen, und zuletzt aus diesen beiden Mitteln das einfache arithmetische
Mittel nehmen müssen.

In Ermangelung einer bestimmten Kenntniss von \(\theta\), ist diese Rechnung
in den drei Hypothesen \(\theta=0\), \(\theta=1\), \(\theta=\infty\) geführt, woraus sich für die
Polhöhe ergeben hat:
\begin{center}
\renewcommand{\arraystretch}{1.18}%
\begin{tabular}{c|c|c|c}
&\(\theta=0\)&\(\theta=1\)&\(\theta=\infty\)\\
\hline
Kreis Ost & \(50^\circ56'5''{,}75\) & \(50^\circ56'5''{,}69\) & \(50^\circ56'5''{,}71\)\\
Kreis West & \hfill \(4{,}62\) & \hfill \(4{,}65\) & \hfill \(4{,}65\)\\
\hline
Polhöhe & \(50\enspace56\enspace5{,}18\) & \(50\enspace56\enspace5{,}17\) & \(50\enspace56\enspace5{,}18\)
\end{tabular}
\end{center}

Man sieht also, dass die Berücksichtigung der strengern Grundsätze das
erste Resultat gar nicht merklich ändert, und dass man sich an die Zahl
\(50^\circ56'5''{,}17\) halten kann.

Bei diesen Rechnungen ist auf die Biegung des Fernrohrs noch keine
Rücksicht genommen. Nach Hrn. \textsc{Hansen}s Angabe ist dieselbe im Horizont
\(=1''{,}00\), und zwar von der beobachteten Zenithdistanz abzuziehen, oder nach
unserer Bezeichnung \(f=-1''{,}00\). Man sieht, dass bei Berücksichtigung dieser
Biegung die Polhöhe aus den nördlich vom Zenith culminirenden Sternen
etwas grösser, aus den südlichen kleiner ausfallen, und, weil jene etwas über-
\par

% Printed 55; principal PDF 62; Google PDF 66.
\par\clearpage
\noindent\makebox[12mm][l]{}\makebox[\dimexpr\linewidth-24mm\relax][c]{\footnotesize\textbf{IV.\quad BREITENBESTIMMUNG DER STERNWARTE SEEBERG.}}\makebox[12mm][r]{\fontsize{13}{16}\selectfont 55}\par\medskip

\noindent wiegen, das Mittelresultat um \(0''{,}02\) vergrössert werden wird. Der zweite Theil
der Biegung, oder die Biegung bei verticaler Stellung, kann, da alle hier
vorkommenden Zenithdistanzen nur klein sind, als eine constante Veränderung
des Collimationsfehlers betrachtet werden, und wird also bei unserm Verfahren
gerade eben so, wie die Theilungsfehler der bei der Bestimmung von jenem
sprechenden Theilstriche, von selbst eliminirt.

Wir haben demnach als Definitivwerth für die Polhöhe aus diesen Beobachtungen
\[
50^\circ56'5''{,}19.
\]

Die erwähnte trigonometrische Verbindung der beiden Sternwarten, nach
den oben angeführten Dimensionen des Erdsphäroids berechnet, gibt den
Breitenunterschied
\[
35'41''{,}86,
\]
\noindent also mit der oben bestimmten Polhöhe von Göttingen die der Seeberger
Sternwarte
\[
=50^\circ56'5''{,}99.
\]

Diese bezieht sich auf den Dreieckspunkt, nemlich das Centrum der Axe
des Mittagsfernrohrs; das Centrum der Axe des Meridiankreises liegt \(1{,}168\)
Toisen, oder im Bogen \(0''{,}07\) südlicher; die Polhöhe des letztern Punktes ist
also, aus Göttingen durch die trigonometrische Verbindung abgeleitet,
\[
=50^\circ56'5''{,}92
\]
\noindent oder \(0''{,}73\) grösser, als aus den astronomischen Beobachtungen.

Für den Längenunterschied folgt übrigens aus der trigonometrischen
Verbindung \(47'9''{,}20\) im Bogen, oder \(3^{\mathrm m}8^{\mathrm s}{,}61\) in Zeit, sehr gut mit unserer
Kenntniss aus astronomischen Beobachtungen übereinstimmend. Endlich folgt aus
jenen Messungen das Azimuth der Dreiecksseite Seeberg---südliches
Meridianzeichen bei Schwabhausen \(4''{,}6\) westlich, welches gleichfalls bei der nicht
unbeträchtlichen Anzahl der Zwischenpunkte, den Verschiedenheiten, die in
den Angaben einiger Winkel der preussischen Messung vorkommen, und der
Ungewissheit, ob der Dreieckspunkt sich genau im Meridian befand, wie eine
gute Übereinstimmung betrachtet werden kann.

\begin{center}\rule{.24\textwidth}{.4pt}\end{center}

% Printed 56 (unfoliated); principal PDF 63; Google PDF 67.
\par\clearpage

\vspace*{.055\textheight}
\begin{center}{\Large ZUSATZ ZU [ART. 20] S. 48.}\end{center}

\textsc{Walbeck}s Bestimmung der Dimensionen des Erdsphäroids befindet sich
in einer kleinen Abhandlung: \textit{De forma et magnitudine telluris, ex dimensis
arcubus meridiani, definiendis}, wovon aber nur die zwei ersten Bogen im Druck
erschienen sind (Åbo, 1819). \textsc{Walbeck} hat die peruanische, die beiden
ostindischen, die französische, englische und die neuere lappländische Gradmessung
dem Calcül unterworfen und ist meines Wissens bisher der einzige, der
dieses Geschäft nach richtigen willkürfreien Grundsätzen ausgeführt hat.
Inzwischen hat er bei jeder einzelnen Gradmessung nur den ganzen Bogen, oder
die an den Endpunkten beobachteten Polhöhen, in Betracht gezogen, ohne
die bei mehrern vorhandenen Zwischenpunkte zu berücksichtigen, und in der
Rechnung ist er bei der ersten Potenz der Abplattung stehen geblieben.

Ich habe deshalb den durch mehrere Arbeiten bereits vortheilhaft
bekannten Hrn. Dr. \textsc{Schmidt} unlängst zu einer neuen Berechnung dieser
sämmtlichen Gradmessungen veranlasst, welche er während des Abdrucks der letzten
Bogen gegenwärtiger Schrift vollendet hat. Er hat dabei sowohl die höhern
Potenzen der Abplattung, als die an allen Zwischenpunkten beobachteten
Polhöhen mit berücksichtigt, auch die hannoversche Gradmessung hinzugezogen,
und, nach dem oben S.~50 angedeuteten Princip, dasjenige Ellipsoid bestimmt,
auf welchem die astronomisch beobachteten Polhöhen, um mit den
geodätischen Messungen in vollkommene Übereinstimmung zu kommen, der
möglich geringsten Abänderung bedürfen, d.\ i. wo die Summe der Quadrate der
hiezu erforderlichen Abänderungen ein Minimum wird. Das Resultat dieser
Rechnung ist:
\par

% Printed 57; principal PDF 64; Google PDF 68.
\par\clearpage
\noindent\makebox[12mm][l]{}\makebox[\dimexpr\linewidth-24mm\relax][c]{\footnotesize\textbf{BESTIMMUNG DES BREITENUNTERSCHIEDES ETC.\quad ZUSATZ.}}\makebox[12mm][r]{\fontsize{13}{16}\selectfont 57}\par\medskip

\begin{center}
\begin{tabular}{@{}l@{\quad}l@{}}
\makebox[.62\linewidth][s]{Abplattung\hfill .\hfill .\hfill .\hfill .\hfill .\hfill .\hfill .\hfill .\hfill .\hfill .\hfill .\hfill .\hfill .}
& \(\dfrac1{298{,}39}\)\\
\makebox[.62\linewidth][s]{Dreihundertsechzigster Theil des Erdmeridians\hfill .}
& \(57010{,}35\) Toisen.
\end{tabular}
\end{center}

Die beobachteten Polhöhen an den 25 Punkten der sieben Gradmessungen
und ihre kleinsten zur vollkommenen Übereinstimmung mit den gefundenen
Erddimensionen erforderlichen Abänderungen stellt folgende Übersicht dar:
\[
\renewcommand{\arraystretch}{1.38}
\begin{array}{@{}l@{\quad}|r@{\enspace}r@{\enspace}r|l@{}}
\multicolumn{5}{c}{\textit{Peruanische Messung.}}\\[.3ex]
\text{Tarqui}&-3^\circ&4'&30''{,}83&+2''{,}05\\
\text{Cotchesqui}&+0&2&37{,}83&-2{,}05\\[.8ex]
\multicolumn{5}{c}{\textit{Erste ostindische Messung.}}\\[.3ex]
\text{Trivandeporum}&+11&44&52{,}59&-0{,}48\\
\text{Paudree}&13&19&49{,}02&+0{,}47\\[.8ex]
\multicolumn{5}{c}{\textit{Zweite ostindische Messung.}}\\[.3ex]
\text{Punnae}&8&9&38{,}39&-1{,}43\\
\text{Putchapolliam}&10&59&48{,}93&-1{,}18\\
\text{Dodagoontah}&12&59&59{,}91&+3{,}37\\
\text{Namthabad}&15&6&0{,}64&-0{,}77\\[.8ex]
\multicolumn{5}{c}{\textit{Französische Messung.}}\\[.3ex]
\text{Formentera}&38&39&56{,}11&+3{,}95\\
\text{Montjouy}&41&21&45{,}45&+2{,}81\\
\text{Barcelona}&41&22&47{,}16&+1{,}07\\
\text{Perpignan}&42&41&58{,}01&-3{,}67\\
\text{Carcassonne}&43&12&54{,}31&-0{,}96\\
\text{Evaux}&46&10&42{,}19&-6{,}14\\
\text{Pantheon}&48&50&48{,}94&-0{,}17\\
\text{Dünkirchen}&51&2&8{,}74&+3{,}12\\[.8ex]
\multicolumn{5}{c}{\textit{Englische Messung.}}\\[.3ex]
\text{Dunnose}&50&37&7{,}81&-1{,}73\\
\text{Greenwich}&51&28&39{,}60&+1{,}00\\
\text{Blenheim}&51&50&27{,}50&+3{,}02\\
\text{Arburyhill}&52&13&27{,}79&+1{,}80\\
\text{Clifton}&53&27&31{,}59&-4{,}07
\end{array}
\]
\noindent\textbf{IX.}\hfill 8\hspace{.12\textwidth}\par

% Printed 58; principal PDF 65; Google PDF 69.
\par\clearpage
\noindent\makebox[12mm][l]{\fontsize{13}{16}\selectfont 58}\makebox[\dimexpr\linewidth-24mm\relax][c]{\footnotesize\textbf{BESTIMMUNG DES BREITENUNTERSCHIEDES ETC.\quad ZUSATZ.}}\makebox[12mm][r]{}\par\medskip

\[
\renewcommand{\arraystretch}{1.38}
\begin{array}{@{}l@{\quad}|r@{\enspace}r@{\enspace}r|l@{}}
\multicolumn{5}{c}{\textit{Hannoversche Messung.}}\\[.3ex]
\text{Göttingen}&51^\circ&31'&47''{,}85&-2''{,}65\\
\text{Altona}&53&32&45{,}27&+2{,}66\\[.8ex]
\multicolumn{5}{c}{\textit{Schwedische Messung.}}\\[.3ex]
\text{Mallörn}&65&31&31{,}06&+1{,}40\\
\text{Pahtavara}&67&8&51{,}41&-1{,}40
\end{array}
\]

Die Zahlen der letzten Columne sind nun keinesweges wie Fehler der
astronomischen Beobachtungen zu betrachten, sondern sie sind die algebraische
Summe dieser Fehler und der Unregelmässigkeiten der Richtung der Verticale.
Wenn man diese Gesammtabweichungen nach denselben Regeln, wie die
zufälligen Fehler, behandelt, so findet sich die mittlere Abweichung \(3''{,}18\), und
damit der mittlere zu befürchtende Fehler
\begin{center}
\begin{tabular}{@{}l@{\quad}l@{}}
\makebox[.64\linewidth][s]{in dem Nenner der Abplattung\hfill .\hfill .\hfill .\hfill .\hfill .\hfill .}
& \(12{,}5\) Einheiten\\
in dem Werthe des dreihundertsechzigsten Theils &\\
\makebox[.64\linewidth][r]{des Erdmeridians} & \(5{,}0\) Toisen.
\end{tabular}
\end{center}

Den sogenannten wahrscheinlichen Fehler mag man also auf 8 Einheiten
bei dem Nenner der Abplattung, und auf 3 Toisen bei dem mittlern
Breitengrade schätzen, und diese Fixirung unserer Begriffe über den Grad der
Genauigkeit, welchen man der Bestimmung der Dimensionen des Erdsphäroids
durch alle bisherigen Breitengradmessungen zuzuschreiben berechtigt ist, hat
man als ein wichtiges Resultat dieser verdienstlichen, an einem andern Orte
ausführlich bekannt zu machenden Arbeit des Hrn. Dr. \textsc{Schmidt} anzusehen.

\begin{center}\rule{.24\textwidth}{.4pt}\end{center}

% Printed 59 (unfoliated); principal PDF 66; Google PDF 70.
\par\clearpage

\vspace*{.075\textheight}
\begin{center}{\LARGE A\,N\,Z\,E\,I\,G\,E.}\end{center}
\vspace{.045\textheight}
\begin{center}
\rule{.52\textwidth}{.4pt}\\[-.2ex]
{\small Göttingische gelehrte Anzeigen.\enspace 1828 Juni 16.}\\[-.2ex]
\rule{.52\textwidth}{.4pt}
\end{center}

Bei Vandenhoeck und Ruprecht: \textit{Bestimmung des Breitenunterschiedes
zwischen den Sternwarten von Göttingen und Altona durch Beobachtungen am}
\textsc{Ramsden}\textit{schen Zenithsector, von} \textsc{Carl Friedrich Gauss}.
1828. 84 Seiten in 4.

Die von dem Verf. dieser Schrift während der verflossenen Jahre im
Königreich Hannover ausgeführten Messungen hatten zunächst den Zweck, die
von dem Hrn. Prof. \textsc{Schumacher} in den dänischen Staaten unternommene
Gradmessung um zwei Grad weiter nach Süden auszudehnen. Eine Dreieckskette
von der südlichen Grenze des Königreichs Hannover bis Hamburg wurde
in den Jahren 1821---1823 vollendet und mit dem dänischen Dreieckssystem
verbunden. Um daraus ein selbstständiges Resultat als Gradmessung ziehen
zu können, war noch die astronomische Bestimmung der Krümmung des ganzen
Meridianbogens erforderlich, welche jedoch, wegen der in den folgenden
Jahren auf höhern Befehl unternommenen Erweiterung des Dreieckssystems
zum Anschluss an die \textsc{Krayenhoff}schen Dreiecke, bis zum Jahr 1827
ausgesetzt bleiben musste. Da die Bekanntmachung der trigonometrischen Arbeiten
aus mehrern hier nicht anzuführenden Gründen einem grössern Werke
vorbehalten bleibt, so hat der Hofr. \textsc{Gauss} Anlass genommen, jenen
astronomischen Theil der Gradmessung jetzt sogleich als ein für sich bestehendes
Werk herauszugeben, damit das Endresultat sofort den übrigen Gradmessungen
beigefügt und benutzt werden könne.

Vortheilhafter liegende Endpunkte, als diejenigen, welche sich hier von
selbst darboten, hätte man sich gar nicht wünschen können. Es sind die
\par\smallskip
\noindent\hfill \(8^{*}\)\hspace{.12\textwidth}\par

% Printed 60; principal PDF 67; Google PDF 71.
\par\clearpage
\noindent\makebox[12mm][l]{\fontsize{13}{16}\selectfont 60}\makebox[\dimexpr\linewidth-24mm\relax][c]{\small\textbf{ANZEIGE.}}\makebox[12mm][r]{}\par\medskip

\noindent Sternwarten von Göttingen und Altona, beide mit trefflichen Instrumenten
ausgerüstet, und beide, durch ein in seiner Art einziges Spiel des Zufalls, so
genau in einem und demselben Meridian liegend, dass man, um einen
Unterschied aufzustellen, bestimmte Plätze in den Sternwarten angeben muss: auf
die Mittelpunkte der Axen der \textsc{Reichenbach}schen Meridiankreise bezogen
liegt nemlich die Altonaer Sternwarte nur \(7\tfrac12\) Toisen westlicher.

Wenn gleich die Ausdehnung des Bogens, welcher dem Stück des
Erdmeridians zwischen diesen beiden Sternwarten am Himmel entspricht, schon
aus den absoluten Polhöhen derselben, zu deren Bestimmung die fortgesetzten
Beobachtungen fortwährend neue Beiträge liefern, sich ergibt, so war es doch
von grosser Wichtigkeit, die Bestimmung des Breitenunterschiedes noch auf
eine andere Art, mit einem und demselben Instrument vom ersten Range, zu
erhalten, und der Hofr. \textsc{Gauss} konnte dazu den trefflichen \textsc{Ramsden}schen
Zenithsector benutzen, welcher zu ähnlichen Operationen bei der englischen
Gradmessung angewandt und bekanntlich von \textsc{Mudge} ausführlich beschrieben ist.

Das Werk zerfällt in vier Abschnitte. Im ersten werden die beobachteten
Sterne nachgewiesen. Es ist der eigenthümliche Character der neuern
Beobachtungskunst, dass bei wichtigen Anlässen die Anzahl der Beobachtungen
sehr vervielfältigt wird, um den Einfluss der Unvollkommenheit der Sinne und
der Instrumente (denn absolut vollkommen \textls[100]{kann} keines sein), so wie der
unvermeidlichen von aussen störenden Ursachen, wenn auch nicht wegzuschaffen
(denn das ist unmöglich), aber doch auf einen sehr kleinen Theil seiner
sonstigen Grösse herabzubringen. Dieser Zweck kann aber nur dadurch erreicht
werden, dass man die einzelnen Fehlerquellen bei einer grossen Menge von
Beobachtungen auf vielfach verschiedene Art ins Spiel treten lässt. Daher
können namentlich die unregelmässigen Theilungsfehler eines nicht
wiederholenden Instruments nur durch Beobachtung einer bedeutenden Anzahl von
Sternen mit ungleichen Declinationen zur bestmöglichen Ausgleichung gebracht
werden, und der Vf. wählte deshalb 43 Sterne in schicklichen Lagen zur
Beobachtung aus, von denen manche noch in keinem Sternverzeichnisse vorkommen.

Der zweite Abschnitt hat die Beobachtungen selbst zum Gegenstande, die
in Göttingen vom Anfang Aprils bis zur Mitte des Mai, und nachher in
Altona während des Junius angestellt wurden, und deren Anzahl sich auf 900
\par

% Printed 61; principal PDF 68; Google PDF 72.
\par\clearpage
\noindent\makebox[12mm][l]{}\parbox[t]{\dimexpr\linewidth-24mm\relax}{\centering\scriptsize\bfseries BESTIMMUNG DES BREITENUNTERSCHIEDES ZWISCHEN GÖTTINGEN UND ALTONA.}\makebox[12mm][r]{\fontsize{13}{16}\selectfont 61}\par\medskip

\noindent beläuft; sie sind alle einzeln aufgeführt, aber nicht in der ursprünglichen
Gestalt des Tagebuchs, was dem Vf. unnöthig schien und den Umfang des
Werks auf das Doppelte vergrössert haben würde, sondern sogleich nach den
Sternen geordnet, nebst ihren Reductionen auf die mittlere Stellung, für den
Anfang des Jahrs 1827.

Der dritte Abschnitt entwickelt die aus den Beobachtungen sich ergebenden
Resultate. Das Hauptresultat ist der Breitenunterschied zwischen den
beiden Sternwarten, welcher, in Beziehung auf die Plätze der
Meridiankreise, nach mehrfachen von verschiedenen Gesichtspunkten ausgehenden
Combinationen auf \(2^\circ0'57''{,}42\) festgesetzt wird, wobei nur eine wahrscheinliche
Unsicherheit von \(0''{,}07\) zurückbleibt. Merkwürdig ist die aus den Beobachtungen
auf das Evidenteste hervorgehende Veränderung des Collimationsfehlers, welcher
bei den Beobachtungen in Göttingen \(3''{,}76\), bei den Beobachtungen in Altona
\(1''{,}39\) betrug, und diese Veränderung auf dem obwohl auf das Vorsichtigste
geleiteten Transport erlitten haben muss. Diese Thatsache ist besonders
wichtig, um das der ältern lappländischen Gradmessung gebührende Zutrauen zu
würdigen, der man mit Recht die Unterlassung des Umwendens des
Zenithsectors an jedem Beobachtungsplatze zum Vorwurfe gemacht hat, und bei der
die Gleichheit des Collimationsfehlers durch das, was darüber neuerlich in
einer sonst sehr schätzbaren Abhandlung vorgebracht ist, keinesweges
befriedigend gerechtfertigt wird. Man findet hier ferner die Beobachtungen des Vfs.
aus den Jahren 1820 und 1824, welche zur Festsetzung der absoluten
Polhöhe der Göttinger Sternwarte dienen; die Vergleichung des aus der
hannoverschen Gradmessung folgenden Resultats für den Breitengrad mit den
Dimensionen des Erdsphäroids, welche \textsc{Walbeck} 1819 aus sämmtlichen zuverlässigen
frühern Gradmessungen abgeleitet hat; die Ansichten des Vfs. von den
Unregelmässigkeiten der Erdfigur überhaupt; endlich die Vergleichung der aus
gegenwärtigen Beobachtungen folgenden Declinationen der Zenithalsterne mit
den \textsc{Bradley}schen und \textsc{Piazzi}schen Bestimmungen. Merkwürdig ist bei einem
der kleinern Sterne (\textsc{Piazzi} 16. 291) eine Differenz von \(10''{,}3\) mit \textsc{Piazzi}s
Catalog, in welcher man nach den vom Vf. beigebrachten Gründen eine
eigene Bewegung nicht verkennen kann.

Der vierte Abschnitt enthält die Bestimmung der Breite der Sternwarte
Seeberg, aus den Beobachtungen, welche der Vorsteher dieser Sternwarte,
\par
\par
\endgroup
